% Generated draft: requires mathematical and source review.
\documentclass[11pt]{article}
\usepackage[margin=25mm]{geometry}
\usepackage{fontspec}
\setmainfont{Latin Modern Roman}
\usepackage{amsmath,amssymb,mathtools,unicode-math}
\setmathfont{Latin Modern Math}
\usepackage{enumitem,needspace,xurl,hyperref}
\setlength{\parindent}{0pt}
\setlength{\parskip}{5pt}
\setlength{\emergencystretch}{3em}
\newcommand{\N}{\mathbb{N}}
\newcommand{\Z}{\mathbb{Z}}
\newcommand{\Q}{\mathbb{Q}}
\newcommand{\R}{\mathbb{R}}
\newcommand{\C}{\mathbb{C}}
\newcommand{\F}{\mathbb{F}}
\newcommand{\A}{\mathbb{A}}
\newcommand{\PP}{\mathbb{P}}
\newcommand{\E}{\mathbb{E}}
\newcommand{\eps}{\varepsilon}
\newcommand{\dd}{\,\mathrm{d}}
\begin{document}
\section*{UnsolvedMath ID: 30000811}
\textbf{Source problem code:} catalogue entry
\subsection{Definitions and mathematical statement}
In the usual finite affine setting, let \(V=\F_p^d\) and \(H\leq\operatorname{GL}(V)\). The group \(G=V\rtimes H\) acts on \(V\) by \((v,h)\cdot x=v+h(x)\). Its action is primitive exactly when \(H\) acts irreducibly on \(V\). Insoluble means that the finite group is not solvable. Under the usual terminology, simply primitive means primitive but not doubly transitive; in this affine setting this requires \(H\) not to be transitive on \(V\setminus\{0\}\).

A primitive permutation group is extremely primitive when each point stabilizer acts primitively on each of its nontrivial orbits away from the fixed point. In the affine setting it is enough to consider the stabilizer \(H=G_0\) and its orbits on \(V\setminus\{0\}\). For a nontrivial such orbit \(Hv\), this is equivalent to \(H_v=\{h\in H:h(v)=v\}\) being a maximal proper subgroup of \(H\). Exact conventions for ``simply primitive'' and nontrivial orbits must be checked against the report.

Let \(\mathcal L_3\) denote the specific list in the report's Theorem 3, which is not supplied in this record. The extracted completeness question is whether every insoluble, simply primitive, extremely primitive finite affine permutation group is permutation-isomorphic to an entry of \(\mathcal L_3\), with no additional exceptions. The record describes the proposed list as groups \(\F_2^d\rtimes H\) associated with sporadic, alternating, or Lie-type socles. Thus completeness would in particular exclude any additional prime, dimension, stabilizer, or affine action beyond that list, not merely place the group in one of three broad families.

There is a crucial qualification about the word socle. For a finite primitive affine group, \(\operatorname{soc}(G)=V\), an elementary abelian regular normal subgroup. It cannot itself be a sporadic, alternating, or Lie-type nonabelian simple group. The cited families must therefore refer to \(H\), or to another group specified in the missing theorem, rather than literally to the affine group's own socle. The cleaned phrase ``whose socle'' does not identify its referent; this cannot be resolved by treating the full affine group as almost simple.
\subsection{Short English statement}
The question asks whether a particular reported list exhausts the nonsolvable affine groups satisfying two strengthened primitivity conditions. The list itself is missing, and the nonabelian socle families must be distinguished from the elementary abelian socle of the full affine group.
\subsection{Sources}
\begin{itemize}
\item Local UnsolvedMath registry, permanent ID \(30000811\), exact source code \texttt{catalogue entry}, proposed year 2007; no individual proposer is supplied.
\item \emph{Permutation Groups (2008)}, Oberwolfach Reports, DOI 10.4171/owr/2007/37, pp. 2173--2236, particularly the original statement's reference to Theorem 3. \url{https://doi.org/10.4171/owr/2007/37}.
\item The supplied \texttt{corrected\_verified} note reports direct confirmation of the affine, insolubility, primitivity, and three socle-family conditions. The literature assessment checked 2026-08-21 nevertheless describes an incomplete conjectural classification fragment and says no complete resolution can be established from the identified literature. The open label is supplied dated triage, not independent verification.
\end{itemize}
\end{document}
